\documentclass[../main2.tex]{subfiles}
\begin{document}

\chapter*{前言（重构版 v2）}
\addcontentsline{toc}{chapter}{前言（重构版 v2）}

\begin{readingnote}
你有没有买过一本讲社会规律的书，读完觉得每个字都对，却什么都预测不了？\\
如果一本书只能解释过去，不能帮你判断接下来会怎样——它算讲清楚了吗？
\end{readingnote}

这本书的第一版叫《社会动力学与公理化体系》。它回答了一个问题：地理、文化、经济、政治、法律，能不能放进同一条推导链？答案是能。但答完之后，我自己不满意。

我不满意三个地方。

第一，它是回溯的。从结果倒推原因，故事很顺。但顺的代价是丢信息——同一个结果可以来自完全不同的原因，倒推没有唯一解。你读着觉得通了，其实只是选了一条顺眼的路。

第二，四个领域写成了并列的四章。文化一章，经济一章，政治一章，法律一章，像四本小书钉在一起。它们之间是什么关系？谁从谁里长出来？没说清。

第三，它没有预测能力。一本讲社会动力的书，说不清“如果当初没有X，今天会怎样”，那就只有后视镜，没有前挡风。

所以有了这一版。这一版是前向的，主线只有一句话：从“想预测未来”这个动作出发，因果推断自然长出来，法律界早把这套方法跑通了，最后把它装到社会运转的四个齿轮上。

和第一版相比，这一版有三处关键改动。我想在开头就挑明，免得你读到一半觉得“怎么变了”。

第一，砍掉了一整层。第一版里有一层叫“演化状态空间”，规定未来由物质、人、相互作用构成。听起来严密，但它不是读者会问的问题。读者问的是“我怎么知道明天会发生什么”，不是“未来是什么东西”。这一版不预设本体论前提，从动机直接起步。

第二，因果从预测里长出来。能预测，就等价于有因果。你能算出“改了A之后会怎样”，说明你知道A是不是原因。反过来也成立。所以本书的目标不是预测趋势，而是预测“社会在不同选择下会走向哪里”。这是干预性预测，只有因果能做。

第三，法律升格为方法的证明。第一版把法律当四个通道之一。这一版把它提前到Part III，单独成一部：法庭必须给出“70\%是他的责任”这种确定数字。这种压力逼着人类练出了最精细的贡献分解工具。Part IV里法律仍作为齿轮之一出现，那是看它在系统里怎么和其他齿轮咬合。两次出现是刻意的两个角度，不是重复。

怎么读这一版：全书六段——导论（为什么算得最准的人预测最离谱）、动机与目标（Part I）、预测与因果（Part II）、法庭上的时光机（Part III）、社会的四大齿轮（Part IV）、回到循环（Part V）。每一章开头都是上一章遗留的问题，每一个新概念都通过“加入一个约束”导出，不凭空引入。

如果你不想读公式，可以先跳过所有“数学建模”盒子。它们是落脚点，不是必经之路。Part I和Part II零公式，公式从Part III起，且前面必有真实案例。

还要交代一句诚实的话：本书给出的是结构，不是预言。它能告诉你一个社会事件“依赖哪几个因素、各占多少”，不能告诉你“明天会发生什么”。方法会在哪里失败，书中有一张失败模式清单——那是本书可信度的一部分。

原书的材料没有被丢弃，全部章节都在新结构里找到了新位置。这是重构，不是重写。它依然在生长，但已经足够诚恳。

\begin{thinkbox}
\textbf{想一想：}
\begin{enumerate}
  \item 回想你最近一次用“因为……所以……”解释一个社会新闻。你说的“因为”是相关，还是因果？你怎么验证？
  \item 如果一本书只能解释过去，不能帮你判断“如果当初没有X，今天会怎样”，你会怎么评价它的用处？反事实（如果没有X会怎样）对你做决定有什么用？
\end{enumerate}
\end{thinkbox}

\end{document}
